\documentclass[aspectratio=169]{beamer}
\input{header}
\coursecode{JKSSB}

\title{Excel Functions}
\subtitle{Lesson 26 --- Core Functions and Error Values}
\author{JKSSB Computer Awareness}
\date{\today}

\begin{document}
\frame{\titlepage}

\begin{frame}{From Formulas to Functions}
  \begin{itemize}
    \item A \key{formula} is an expression that starts with \texttt{=} and
      calculates a value, for example \texttt{=A1+A2}.
    \item A \key{function} is a ready-made, named operation such as
      \texttt{SUM} or \texttt{AVERAGE} that performs the calculation for you.
    \item A function is used \key{inside} a formula; its name is followed by
      arguments in parentheses.
    \item Using a function is faster and less error-prone than typing every
      cell into the formula by hand.
  \end{itemize}
  \begin{block}{The one idea}
    A formula is the whole calculation; a function is a built-in tool that
    performs one part of it.
  \end{block}
\end{frame}

\begin{frame}{Function vs Formula}
  \begin{center}
    \begin{tabular}{p{0.20\textwidth}p{0.32\textwidth}p{0.34\textwidth}}
      \toprule
      \textbf{Point} & \textbf{Formula} & \textbf{Function} \\
      \midrule
      What it is & Any calculation starting with \texttt{=} & A predefined named operation \\
      Example & \texttt{=A1+A2} & \texttt{=SUM(A1:A10)} \\
      Contains a function? & Not always & Always (by definition) \\
      \bottomrule
    \end{tabular}
  \end{center}
  \vspace{0.25cm}
  \begin{alertblock}{Trap alert}
    Every formula must begin with \texttt{=}. Without it, Excel treats the
    entry as plain text. (TRAP-30)
  \end{alertblock}
\end{frame}

\begin{frame}{Anatomy of a Function}
  \begin{itemize}
    \item A function has a \key{name} and \key{arguments} inside parentheses:
      \texttt{=FUNCTION(argument1, argument2)}.
    \item The \key{name} tells Excel what to do --- \texttt{SUM},
      \texttt{MAX}.
    \item The \key{arguments} tell it which cells or values to use, such as
      the range \texttt{A1:A10}.
    \item A \key{range} is written with a colon: \texttt{A1:A10} means the
      cells from A1 down to A10.
  \end{itemize}
  \begin{exampleblock}{Worked example}
    \texttt{=SUM(B2:B6)} adds the values in cells B2, B3, B4, B5, and B6.
  \end{exampleblock}
\end{frame}

\begin{frame}{SUM: Adding Values}
  \begin{itemize}
    \item \texttt{SUM} returns the \key{total} of the numbers in its
      arguments.
    \item Example: \texttt{=SUM(A1:A10)} adds all the numbers in the range
      A1:A10.
    \item Text and blank cells in the range are \muted{ignored}; only numbers
      are added.
    \item Several ranges can be listed: \texttt{=SUM(A1:A5, C1:C5)}.
  \end{itemize}
  \begin{block}{Remember}
    SUM means addition --- it is the function behind the AutoSum button.
  \end{block}
\end{frame}

\begin{frame}{AVERAGE: The Mean}
  \begin{itemize}
    \item \texttt{AVERAGE} returns the \key{arithmetic mean} of the numbers
      in a range.
    \item Example: \texttt{=AVERAGE(A1:A10)} adds the numbers and divides by
      how many numbers there are.
    \item Blank cells and text are \key{not counted} in the divisor.
    \item Trap: it divides by the number of \key{numeric} cells, not by the
      total number of cells.
  \end{itemize}
  \begin{exampleblock}{Worked example}
    If A1:A5 holds 10, 20, 30 and two blanks, \texttt{=AVERAGE(A1:A5)} is
    $(10+20+30)/3 = 20$, not $60/5$.
  \end{exampleblock}
\end{frame}

\begin{frame}{COUNT vs COUNTA: The Near-Neighbour}
  \begin{center}
    \begin{tabular}{p{0.20\textwidth}p{0.32\textwidth}p{0.34\textwidth}}
      \toprule
      \textbf{Point} & \textbf{COUNT} & \textbf{COUNTA} \\
      \midrule
      Counts & Only cells holding numbers & All non-empty cells \\
      Text cells & Ignored & Counted \\
      Blank cells & Ignored & Ignored \\
      Example & \texttt{=COUNT(A1:A10)} & \texttt{=COUNTA(A1:A10)} \\
      \bottomrule
    \end{tabular}
  \end{center}
  \vspace{0.2cm}
  \muted{COUNT counts numbers; COUNTA counts anything that is not empty.}
\end{frame}

\begin{frame}{COUNT vs COUNTA in Action}
  \begin{itemize}
    \item Suppose A1:A10 holds \key{6 numbers}, \key{2 text} labels, and
      \key{2 blank} cells.
    \item \texttt{=COUNT(A1:A10)} $=$ \key{6} (numbers only).
    \item \texttt{=COUNTA(A1:A10)} $=$ \key{8} (6 numbers $+$ 2 text
      labels).
    \item Both ignore the blanks; only COUNTA counts the text.
  \end{itemize}
  \begin{alertblock}{Trap alert}
    COUNT counts \key{numbers only}; COUNTA counts \key{all non-empty} cells.
    Do not swap this near-neighbour pair.
  \end{alertblock}
\end{frame}

\begin{frame}{MAX and MIN}
  \begin{itemize}
    \item \texttt{MAX} returns the \key{largest} number in a range.
    \item \texttt{MIN} returns the \key{smallest} number in a range.
    \item Examples: \texttt{=MAX(C1:C10)} and \texttt{=MIN(C1:C10)}.
    \item Blank cells and text are ignored, exactly as in SUM and AVERAGE.
  \end{itemize}
  \begin{block}{Pair them}
    MAX $\rightarrow$ biggest; MIN $\rightarrow$ smallest. Both work on
    numbers only.
  \end{block}
\end{frame}

\begin{frame}{IF: Making a Decision}
  \begin{itemize}
    \item \texttt{IF} tests a condition and returns one of two values.
    \item Syntax: \texttt{=IF(logical\_test, value\_if\_true,
      value\_if\_false)}.
    \item Example: \texttt{=IF(A1>50,"Pass","Fail")} returns
      \texttt{Pass} when A1 is more than 50.
    \item The first argument must be a \key{logical test} that is either true
      or false.
  \end{itemize}
  \begin{exampleblock}{Worked example}
    \texttt{=IF(B2>=40,"Pass","Fail")} shows \texttt{Pass} for a mark of 40
    or above and \texttt{Fail} below it.
  \end{exampleblock}
\end{frame}

\begin{frame}{Nested IF and Blank Input}
  \begin{itemize}
    \item A \key{nested IF} places an IF inside the value-if-false part, to
      test more than one condition.
    \item Example: \texttt{=IF(A1>50,"High",IF(A1>25,"Medium","Low"))}.
    \item Excel checks the conditions \key{in order}; the first true test
      wins.
    \item In a numeric comparison Excel reads a \key{blank} cell as
      \key{0}, so a blank fails the tests and the final value is returned.
  \end{itemize}
  \begin{alertblock}{PYQ alert}
    With A1 blank, the formula above returns \key{Low}: the blank is read as
    0. (PYQ-14, JA 2026 Q74)
  \end{alertblock}
\end{frame}

\begin{frame}{Other Basic Functions Worth Knowing}
  \begin{center}
    \begin{tabular}{ll}
      \toprule
      \textbf{Function} & \textbf{What it does} \\
      \midrule
      \texttt{COUNTIF} & Counts cells meeting one condition \\
      \texttt{SUMIF} & Adds cells meeting one condition \\
      \texttt{MEDIAN} & Returns the middle value of the numbers \\
      \texttt{ROUND} & Rounds a number to a set number of digits \\
      \texttt{LEN} & Counts the characters in a text value \\
      \texttt{TODAY} & Returns today's date \\
      \bottomrule
    \end{tabular}
  \end{center}
  \vspace{0.2cm}
  \muted{At JKSSB depth, know what each one returns and that all are written
  inside a formula.}
\end{frame}

\begin{frame}{AutoSum}
  \begin{itemize}
    \item \key{AutoSum} is a button that inserts a \texttt{SUM} formula for
      the cells beside the active cell.
    \item By default it writes \texttt{=SUM(...)} over the range Excel
      guesses you want.
    \item Keyboard shortcut: \key{Alt + =} inserts AutoSum.
    \item Its drop-down also offers \texttt{Average}, \texttt{Count Numbers},
      \texttt{Max}, and \texttt{Min}.
  \end{itemize}
  \begin{alertblock}{Shortcut alert}
    AutoSum is \key{Alt + =}, not a Ctrl combination or a function key. This
    is a common wrong option.
  \end{alertblock}
\end{frame}

\begin{frame}{Excel Error Values (1)}
  \begin{center}
    \begin{tabular}{p{0.22\textwidth}p{0.60\textwidth}}
      \toprule
      \textbf{Error} & \textbf{Meaning} \\
      \midrule
      \texttt{\#DIV/0!} & A number is divided by zero or by a blank cell \\
      \texttt{\#VALUE!} & The wrong type of argument --- text where a number is needed \\
      \texttt{\#NAME?} & Excel does not recognise a name --- often a misspelt function name \\
      \bottomrule
    \end{tabular}
  \end{center}
  \vspace{0.2cm}
  \muted{Read the error as a clue to what Excel could not do.}
\end{frame}

\begin{frame}{Excel Error Values (2)}
  \begin{center}
    \begin{tabular}{p{0.22\textwidth}p{0.60\textwidth}}
      \toprule
      \textbf{Error} & \textbf{Meaning} \\
      \midrule
      \texttt{\#REF!} & A cell reference is not valid --- the referenced cell was deleted \\
      \texttt{\#N/A} & The value is not available, as when a lookup finds no match \\
      \texttt{\#NUM!} & A number in the formula is not valid for the operation \\
      \bottomrule
    \end{tabular}
  \end{center}
  \vspace{0.2cm}
  \muted{\texttt{\#\#\#\#\#} is not an error value --- it only means the column is too narrow to show the number.}
\end{frame}

\begin{frame}{Trap Alert: What to Keep Straight}
  \begin{itemize}
    \item A formula must start with \texttt{=}; without it the entry is text.
      (TRAP-30)
    \item \texttt{COUNT} counts \key{numbers}; \texttt{COUNTA} counts all
      \key{non-empty} cells.
    \item \texttt{AVERAGE} divides by the number of numeric cells, not by all
      the cells.
    \item A \key{function} is a named built-in operation used \key{inside} a
      formula.
    \item \texttt{\#DIV/0!} is division by zero; \texttt{\#VALUE!} is a wrong
      data type; \texttt{\#NAME?} is an unrecognised name.
  \end{itemize}
\end{frame}

\begin{frame}{Office Desktop Example}
  \begin{itemize}
    \item A marks sheet lists student names in A2:A31 and marks in B2:B31.
    \item \texttt{=COUNTA(A2:A31)} counts how many students were entered.
    \item \texttt{=COUNT(B2:B31)} counts how many have a numeric mark.
    \item \texttt{=AVERAGE(B2:B31)} gives the class average;
      \texttt{=MAX(B2:B31)} gives the top mark.
  \end{itemize}
  \begin{exampleblock}{Read the two counts together}
    If COUNTA gives 30 but COUNT gives 28, two students have a non-numeric
    entry in the marks column --- a data-entry problem to fix.
  \end{exampleblock}
\end{frame}

\begin{frame}{Lesson 26: Recall}
  \begin{itemize}
    \item A formula starts with \texttt{=}; a function such as \texttt{SUM}
      or \texttt{IF} is a named operation used inside a formula.
    \item \texttt{SUM} adds, \texttt{AVERAGE} means, \texttt{MAX}/\texttt{MIN}
      give the largest/smallest number.
    \item \texttt{COUNT} counts numbers; \texttt{COUNTA} counts all
      non-empty cells --- the key near-neighbour pair.
    \item \texttt{IF(logical\_test, value\_if\_true, value\_if\_false)};
      nested IF checks conditions in order, and a blank is read as 0 in a
      numeric test (PYQ-14).
    \item AutoSum inserts \texttt{=SUM(...)}; the shortcut is \texttt{Alt + =}.
    \item Error values: \texttt{\#DIV/0!}, \texttt{\#VALUE!},
      \texttt{\#NAME?}, \texttt{\#REF!}, \texttt{\#N/A}.
  \end{itemize}
\end{frame}
\end{document}
